\documentclass[10pt]{extarticle}
\usepackage[a4paper,margin=1.5cm]{geometry}
\usepackage[T1]{fontenc}
\usepackage{amsmath,amssymb}
\usepackage{graphicx}
\usepackage[export]{adjustbox}
\usepackage{multicol}
\usepackage{enumitem}
\usepackage{booktabs,array,tabularx}
\setlength{\columnsep}{0.7cm}
\setlength{\parindent}{0pt}
\setlength{\parskip}{3pt}
\pagestyle{empty}
\makeatletter
% figure registry (figdefs.tex is generated by the extract step)
\newcommand{\qfigdef}[4]{\@namedef{qf@file@#1}{#2}\@namedef{qf@w@#1}{#3}\@namedef{qf@wide@#1}{#4}}
\newcommand{\qf@check}[1]{\@ifundefined{qf@file@#1}{\PackageError{score2}{Unknown figure id #1}{}}{}}
\newcommand{\qf@img}[4][]{\includegraphics[width=#2\dimexpr\@nameuse{qf@w@#4}pt\relax,max width=#3,max height=0.45\textheight,#1]{\@nameuse{qf@file@#4}}}
% Full-width questions: finalize wraps every question that contains a wide figure in \qfullbegin...\qfullend,
% so the WHOLE question (stem, figures, options) is set across the page in reading order, then the two columns
% resume. The question's head - its number, stem and first figure (\qfullsplit marks the end of the head) -
% is kept on one page, so a number is never left at a page foot with its figure overleaf; the rest may break
% across pages normally, so no big gaps are left. Without a \qfullsplit, the whole question is kept together
% if it is under 45% of a page. Never write these yourself.
\newif\ifqf@full
\newif\ifqf@boxed
\newsavebox{\qf@box}
\newcommand{\qf@open}{\global\qf@fulltrue\par\global\qf@boxedtrue
  \setbox\qf@box\vbox\bgroup\hsize\textwidth\linewidth\textwidth}
\newcommand{\qf@emit}[1]{\par\egroup\global\qf@boxedfalse
  \ifdim\dimexpr\ht\qf@box+\dp\qf@box\relax<#1\textheight\noindent\usebox{\qf@box}\par\else\unvbox\qf@box\fi}
\newcommand{\qfullsplit}{\ifqf@boxed\qf@emit{0.85}\fi}
\newcommand{\qf@close}{\ifqf@boxed\qf@emit{0.45}\fi\global\qf@fullfalse}
\newcommand{\qfullbegin}{\end{multicols}\qf@open}
\newcommand{\qfullend}{\qf@close\begin{multicols}{2}}
\newcommand{\qfullstart}{\qf@open}
\newcommand{\qfulldone}{\qf@close}
% \qfig[scale]{id}: block figure, centred, directly where it is written (size set by the tool).
% Fallback if a wide figure is not inside a full-width question: break out of the columns in place - no floats.
\newcommand{\qfig}[2][1]{\qf@check{#2}%
  \ifqf@full\qf@blk{#1}{#2}\else
  \ifnum\@nameuse{qf@wide@#2}=1 \expandafter\@firstoftwo\else\expandafter\@secondoftwo\fi
  {\end{multicols}\begin{center}\qf@img{#1}{\textwidth}{#2}\end{center}\begin{multicols}{2}}%
  {\qf@blk{#1}{#2}}\fi}
\newcommand{\qf@blk}[2]{\par\smallskip{\centering\qf@img{#1}{\linewidth}{#2}\par}\smallskip}
% \qfigi[scale]{id}: inline figure (inside option lists / tables / minipages); never breaks columns.
\newcommand{\qfigi}[2][1]{\qf@check{#2}\qf@img[valign=c]{#1}{\linewidth}{#2}}
\makeatother
% structure macros - lint reads these to check completeness and figure placement
\newcommand{\papersec}[1]{\par\bigskip{\centering\Large\bfseries PAPER-#1\par}\medskip}
\newcommand{\partsec}[2]{\par\medskip\noindent{\large\bfseries PART-#1: #2}\par\smallskip}
\newcommand{\sect}[1]{\par\smallskip\noindent\textbf{#1}\par}
\newcommand{\qn}[1]{\par\smallskip\noindent\textbf{#1.}~}
\newcommand{\paraq}[2]{\par\smallskip\noindent\textbf{Paragraph for Questions #1 and #2}\par}
\newenvironment{qoptions}{\begin{itemize}[leftmargin=1.8em,itemsep=1pt,topsep=2pt,label={}]}{\end{itemize}}
\newcommand{\opt}[2]{\item[(#1)] #2}
% figure registry (generated by score2_tools.py finalize)
\qfigdef{p1_phy_s2_q05_1}{TNAS - Test 02 (07 May 2023)-img/p1_phy_s2_q05_1.png}{77.4}{0}
\qfigdef{p1_phy_s2_q06_1}{TNAS - Test 02 (07 May 2023)-img/p1_phy_s2_q06_1.png}{145.2}{0}
\qfigdef{p1_phy_s2_q08_1}{TNAS - Test 02 (07 May 2023)-img/p1_phy_s2_q08_1.png}{111.1}{0}
\qfigdef{p1_phy_s2_q09_1}{TNAS - Test 02 (07 May 2023)-img/p1_phy_s2_q09_1.png}{107.8}{0}
\qfigdef{p1_phy_s2_q10_1}{TNAS - Test 02 (07 May 2023)-img/p1_phy_s2_q10_1.png}{121.7}{0}
\qfigdef{p1_phy_s3_para3_4_1}{TNAS - Test 02 (07 May 2023)-img/p1_phy_s3_para3_4_1.png}{108.5}{0}
\qfigdef{p1_phy_s3_para5_6_1}{TNAS - Test 02 (07 May 2023)-img/p1_phy_s3_para5_6_1.png}{76.8}{0}
\qfigdef{p1_phy_s4_q07_1}{TNAS - Test 02 (07 May 2023)-img/p1_phy_s4_q07_1.png}{113.1}{0}
\qfigdef{p1_phy_s4_q08_1}{TNAS - Test 02 (07 May 2023)-img/p1_phy_s4_q08_1.png}{113.7}{0}
\qfigdef{p1_phy_s4_q09_1}{TNAS - Test 02 (07 May 2023)-img/p1_phy_s4_q09_1.png}{89.6}{0}
\qfigdef{p1_chem_s1_q02_A}{TNAS - Test 02 (07 May 2023)-img/p1_chem_s1_q02_A.png}{58.2}{0}
\qfigdef{p1_chem_s1_q02_B}{TNAS - Test 02 (07 May 2023)-img/p1_chem_s1_q02_B.png}{55.6}{0}
\qfigdef{p1_chem_s1_q02_C}{TNAS - Test 02 (07 May 2023)-img/p1_chem_s1_q02_C.png}{66.9}{0}
\qfigdef{p1_chem_s1_q03_1}{TNAS - Test 02 (07 May 2023)-img/p1_chem_s1_q03_1.png}{101.5}{0}
\qfigdef{p1_chem_s1_q03_A}{TNAS - Test 02 (07 May 2023)-img/p1_chem_s1_q03_A.png}{65.9}{0}
\qfigdef{p1_chem_s1_q03_B}{TNAS - Test 02 (07 May 2023)-img/p1_chem_s1_q03_B.png}{39.9}{0}
\qfigdef{p1_chem_s1_q03_C}{TNAS - Test 02 (07 May 2023)-img/p1_chem_s1_q03_C.png}{39.6}{0}
\qfigdef{p1_chem_s1_q03_D}{TNAS - Test 02 (07 May 2023)-img/p1_chem_s1_q03_D.png}{39.9}{0}
\qfigdef{p1_chem_s3_para3_4_1}{TNAS - Test 02 (07 May 2023)-img/p1_chem_s3_para3_4_1.png}{60.8}{0}
\qfigdef{p1_chem_s4_q08_1}{TNAS - Test 02 (07 May 2023)-img/p1_chem_s4_q08_1.png}{148.0}{0}
\qfigdef{p1_chem_s4_q09_1}{TNAS - Test 02 (07 May 2023)-img/p1_chem_s4_q09_1.png}{75.2}{0}
\qfigdef{p2_phy_s1_q01_1}{TNAS - Test 02 (07 May 2023)-img/p2_phy_s1_q01_1.png}{108.3}{0}
\qfigdef{p2_phy_s1_q01_2}{TNAS - Test 02 (07 May 2023)-img/p2_phy_s1_q01_2.png}{108.3}{0}
\qfigdef{p2_phy_s1_q02_1}{TNAS - Test 02 (07 May 2023)-img/p2_phy_s1_q02_1.png}{97.7}{0}
\qfigdef{p2_phy_s1_q03_1}{TNAS - Test 02 (07 May 2023)-img/p2_phy_s1_q03_1.png}{71.5}{0}
\qfigdef{p2_phy_s1_q04_1}{TNAS - Test 02 (07 May 2023)-img/p2_phy_s1_q04_1.png}{105.1}{0}
\qfigdef{p2_phy_s1_q04_2}{TNAS - Test 02 (07 May 2023)-img/p2_phy_s1_q04_2.png}{78.9}{0}
\qfigdef{p2_phy_s1_q05_1}{TNAS - Test 02 (07 May 2023)-img/p2_phy_s1_q05_1.png}{121.5}{0}
\qfigdef{p2_phy_s1_q06_1}{TNAS - Test 02 (07 May 2023)-img/p2_phy_s1_q06_1.png}{96.8}{0}
\qfigdef{p2_phy_s2_para9_10_1}{TNAS - Test 02 (07 May 2023)-img/p2_phy_s2_para9_10_1.png}{116.5}{0}
\qfigdef{p2_phy_s3_para1_2_1}{TNAS - Test 02 (07 May 2023)-img/p2_phy_s3_para1_2_1.png}{87.0}{0}
\qfigdef{p2_phy_s3_para3_4_1}{TNAS - Test 02 (07 May 2023)-img/p2_phy_s3_para3_4_1.png}{66.3}{0}
\qfigdef{p2_phy_s3_para5_6_1}{TNAS - Test 02 (07 May 2023)-img/p2_phy_s3_para5_6_1.png}{76.5}{0}
\qfigdef{p2_phy_s3_para5_6_2}{TNAS - Test 02 (07 May 2023)-img/p2_phy_s3_para5_6_2.png}{86.7}{0}
\qfigdef{p2_phy_s3_para5_6_3}{TNAS - Test 02 (07 May 2023)-img/p2_phy_s3_para5_6_3.png}{102.6}{0}
\qfigdef{p2_phy_s4_q08_1}{TNAS - Test 02 (07 May 2023)-img/p2_phy_s4_q08_1.png}{89.7}{0}
\qfigdef{p2_phy_s4_q09_1}{TNAS - Test 02 (07 May 2023)-img/p2_phy_s4_q09_1.png}{97.9}{0}
\qfigdef{p2_chem_s1_q01_B}{TNAS - Test 02 (07 May 2023)-img/p2_chem_s1_q01_B.png}{60.8}{0}
\qfigdef{p2_chem_s1_q01_C}{TNAS - Test 02 (07 May 2023)-img/p2_chem_s1_q01_C.png}{42.4}{0}
\qfigdef{p2_chem_s1_q01_D}{TNAS - Test 02 (07 May 2023)-img/p2_chem_s1_q01_D.png}{63.5}{0}
\qfigdef{p2_chem_s1_q01_4}{TNAS - Test 02 (07 May 2023)-img/p2_chem_s1_q01_4.png}{113.1}{0}
\qfigdef{p2_chem_s1_q02_A}{TNAS - Test 02 (07 May 2023)-img/p2_chem_s1_q02_A.png}{44.3}{0}
\qfigdef{p2_chem_s1_q02_B}{TNAS - Test 02 (07 May 2023)-img/p2_chem_s1_q02_B.png}{44.8}{0}
\qfigdef{p2_chem_s1_q02_C}{TNAS - Test 02 (07 May 2023)-img/p2_chem_s1_q02_C.png}{51.1}{0}
\qfigdef{p2_chem_s1_q02_D}{TNAS - Test 02 (07 May 2023)-img/p2_chem_s1_q02_D.png}{51.4}{0}
\qfigdef{p2_chem_s1_q03_1}{TNAS - Test 02 (07 May 2023)-img/p2_chem_s1_q03_1.png}{118.7}{0}
\qfigdef{p2_chem_s3_para1_2_1}{TNAS - Test 02 (07 May 2023)-img/p2_chem_s3_para1_2_1.png}{151.1}{0}
\qfigdef{p2_chem_s3_q01_1}{TNAS - Test 02 (07 May 2023)-img/p2_chem_s3_q01_1.png}{75.7}{0}
\qfigdef{p2_chem_s3_q02_1}{TNAS - Test 02 (07 May 2023)-img/p2_chem_s3_q02_1.png}{95.1}{0}
\qfigdef{p2_chem_s3_para3_4_1}{TNAS - Test 02 (07 May 2023)-img/p2_chem_s3_para3_4_1.png}{124.1}{0}
\qfigdef{p2_chem_s3_q04_1}{TNAS - Test 02 (07 May 2023)-img/p2_chem_s3_q04_1.png}{225.3}{0}
\qfigdef{p2_chem_s4_q07_1}{TNAS - Test 02 (07 May 2023)-img/p2_chem_s4_q07_1.png}{115.8}{0}
\raggedcolumns
\makeatletter\@beginparpenalty=10000\makeatother
\newcolumntype{L}[1]{>{\raggedright\arraybackslash}p{#1}}
\begin{document}
\begin{center}{\Large\bfseries TNAS --- Test 02 (07 May 2023)}\end{center}
\begin{multicols}{2}
\raggedcolumns
\papersec{1}
\partsec{1}{PHYSICS}
\sect{SECTION-I(i)}
\qn{1} A current $I=20$ ampere flows through a wire shaped in the form of an infinite parabola of latus rectum $4a$ ($a=1$ mm). The magnetic field at the focus of the parabola is
\begin{qoptions}
\opt{A}{$\dfrac{\pi}{2}\times10^{-3}$ T}
\opt{B}{$\pi\times10^{-3}$ T}
\opt{C}{$2\pi\times10^{-3}$ T}
\opt{D}{$4\pi\times10^{-3}$ T}
\end{qoptions}
\qn{2} A beam of light of width $t$ is incident on an air-water boundary at an angle of incidence $45^\circ$. The width of the beam in water is (Refractive index of water $=\mu$)
\begin{qoptions}
\opt{A}{$(\mu-1)t$}
\opt{B}{$\dfrac{\sqrt{2\mu^2-1}}{\sqrt2\,\mu}\cdot t$}
\opt{C}{$\dfrac{\sqrt{2\mu^2-1}}{\mu}\cdot t$}
\opt{D}{$\dfrac{\sqrt{\mu^2-1}}{\mu}\cdot t$}
\end{qoptions}
\qn{3} Consider an isolated sphere of radius $b$, made of a material with permittivity $\varepsilon_0$ and electrical conductivity $\gamma$, whose volume charge density varies with time $t$ and is given as
\[\rho=\begin{cases}\rho_0e^{-t/\tau} & \text{if } 0\leqslant r<a\\ 0 & \text{if } a<r<\infty\end{cases}\]
Where $\tau=\dfrac{\varepsilon_0}{\gamma}$ and $0<a<b$. The spherical layer between radii $a$ and $b$ is neutral. Then choose the \textbf{incorrect} option.
\begin{qoptions}
\opt{A}{The surface charge density at outer surface of sphere is $\dfrac{a^3\rho_0}{8b^2}$ at $t=\tau\,\ell n2$.}
\opt{B}{The total current in the region $0<r<a$ increases with the radius.}
\opt{C}{The electric field in the region $b<r<\infty$ is independent of time.}
\opt{D}{The total current in the region $a<r<b$ is independent of radius.}
\end{qoptions}
\qn{4} An electron accelerated by a potential difference $V=3.6$ volt, first enters a region of uniform electric field of a parallel-plate capacitor whose plates extend over a length $\ell=6$ cm in the direction of initial velocity. The electric field is normal to the direction of initial velocity and its strength varies with time as $E=a\times t$ where $a=3200$ Vm$^{-1}$s$^{-1}$. Then the electron enters a region of uniform magnetic field of induction $B=\pi\times10^{-9}$ T. Direction of magnetic field is same as that of the electric field. Calculate the pitch (in mm) of helical path traced by the electron in the magnetic field.\\
(Mass of electron, $m=9\times10^{-31}$ kg, $e=1.6\times10^{-19}$ C). [Neglect the effect of induced magnetic field]
\begin{qoptions}
\opt{A}{$4$}
\opt{B}{$5$}
\opt{C}{$9$}
\opt{D}{$3$}
\end{qoptions}
\sect{SECTION-I(ii)}
\qn{5} An infinite line charge carrying a uniform line charge of density $\lambda$ is placed along the body diagonal AG of a cube of side $a$ as shown. Then
\qfig{p1_phy_s2_q05_1}
\begin{qoptions}
\opt{A}{The flux of the electric field of this charge through one of the faces of the cube is $\dfrac{\lambda a}{2\sqrt3\,\varepsilon_0}$.}
\opt{B}{The flux of the electric field of this charge through one of the faces of the cube is $\dfrac{\sqrt3\lambda a}{2\varepsilon_0}$.}
\opt{C}{The flux of the electric field of this charge through one of the faces of the cube is $\dfrac{\lambda a}{4\sqrt3\,\varepsilon_0}$.}
\opt{D}{The flux of the electric field through the face ABFE is same as that through BCGF.}
\end{qoptions}
\qn{6} A very long wire abcdef carrying current has the shape as shown in the figure. The magnetic field at point P in terms of unit vectors.
\qfig{p1_phy_s2_q06_1}
\begin{qoptions}
\opt{A}{Due to wire ab is $\dfrac{\mu_0I}{4\pi a}\left[1-\dfrac{1}{\sqrt5}\right]\hat k$}
\opt{B}{Due to wire bc is $\dfrac{-\mu_0I}{\pi a\sqrt5}\hat k$}
\opt{C}{Due to wire cd is $\dfrac{-\mu_0I}{\pi a\sqrt5}\hat j$}
\opt{D}{Due to wire ef is $\dfrac{-\mu_0I}{2\pi a}\left[1+\dfrac{2}{\sqrt5}\right]\hat j$}
\end{qoptions}
\qn{7} Two infinite long thin wire have linear charge density $-\lambda$ and $+\lambda$. They are parallel to $z$-axis passing through $x$-axis at point $x=-a$ and $x=+a$ respectively. The equipotential surface having potential $\dfrac{\lambda\ell n(2)}{4\pi\varepsilon_0}$ is
\begin{qoptions}
\opt{A}{a sphere}
\opt{B}{A cylinder}
\opt{C}{having radius of $\sqrt2a$}
\opt{D}{having radius $2\sqrt2a$}
\end{qoptions}
\qn{8} A positive point charge having charge $Q$ is placed at a distance $a/2$ from centre of an uncharged spherical conductor of inner radius $a$ and outer radius $2a$. Another positive point charge having charge $Q$ is placed at distance $4a$ from centre of a conducting sphere. Then which of the following is correct ?
\qfig{p1_phy_s2_q08_1}
\begin{qoptions}
\opt{A}{The surface charge density on inner and outer surface of sphere would be non uniform.}
\opt{B}{The electrostatic potential on outer surface of conductor is $\dfrac{3KQ}{2a}$.}
\opt{C}{Electric potential at centre of sphere is $\dfrac{7KQ}{4a}$.}
\opt{D}{Electric field inside cavity is non-uniform.}
\end{qoptions}
\qn{9} Consider a satellite S which is just above the Brazil and rotating on equatorial plane of earth. The sense of rotation of earth and satellite are similar means both are rotating from west to east. The time period of satellite is 6 hours. At $t=0$ satellite is just above the Brazil. Now choose the correct option(s)
\qfig{p1_phy_s2_q09_1}
\begin{qoptions}
\opt{A}{The time after which the satellite will be above the Brazil again is 8 hours.}
\opt{B}{The above satellite may be a geostationary satellite}
\opt{C}{The height of geostationary satellite above the sea level is 35,800 km approximately.}
\opt{D}{If the distance from the centre of earth to the satellite S becomes $4r_0$, then the time after which the satellite will be above the Brazil again is 24 hours.}
\end{qoptions}
\qn{10} A capacitor with air as dielectric has capacity $C$ as shown in figure is completely filled with dielectric slab of mass $M$ and is connected to a block of mass $M$. The system is released form rest at time $t=0$ when the capacitor is completely filled with dielectric. Then: (Dielectric constant of slab $=K$)
\qfig{p1_phy_s2_q10_1}
\begin{qoptions}
\opt{A}{Charge on the plate of the capacitor when it is half filled with air is $\dfrac{(K+1)CV}{2}$}
\opt{B}{The kinetic energy of the block when the capacitor is half filled with dielectric is $\dfrac{mg\ell}{4}-\dfrac{CV^2(K-1)}{8}$}
\opt{C}{The potential difference between the plates when the capacitor is half filled with air is $V$.}
\opt{D}{The energy stored in the capacitor when it is half filled with dielectric is $K/2\,CV^2$.}
\end{qoptions}
\sect{SECTION-II(i)}
\paraq{1}{2}
\textbf{Rear view mirrors}\\
Rear view mirrors in automobiles are generally convex. Suppose a car A moves with a constant speed of 40.0 kilometer per hour on a straight level road and is followed by another car B moving with the constant speed 60.0 kilometer per hour. At a given instant of time, we denote :\\
$x$ : distance of the car B from the mirror of car A,\\
$y$ : distance of the car B from A as seen by the driver of A in the mirror,\\
$v_x$ : speed of approach of B relative to A and\\
$v_y$ : speed of approach of B as seen in the mirror of A.
\qn{1} If expression for $v_y$ in terms of $x$, $v_x$ and the radius of curvature $R$ of the convex mirror used as the rear view mirrors in the car A is $v_y=\dfrac{R^2}{n(x+0.5R)^2}v_x$. Find $n$.
\qn{2} If $R=2.0$ m, what is the speed of approach of B in kilometer per hour as seen by the driver of A in the mirror for $x=2.0$ m ?
\paraq{3}{4}
In the given circuit diagram, both capacitors are initially uncharged. The capacitance $C_1=2\,\mu$F and $C_2=4\,\mu$F, emf of battery A and B are 2V and 4V respectively.
\qfig{p1_phy_s3_para3_4_1}
\qn{3} On closing switch $S_1$ with $S_2$ open work done by battery A is (in \underline{\hspace{1.2em}} $\mu$J)
\qn{4} On closing switch $S_2$ with $S_1$ open work done by battery B is (in $\mu$J)
\paraq{5}{6}
An isosceles right angle tunnel ABO is made inside the earth from surface of earth to its centre as shown in the figure. The distance between the points A and B and that between the points B and O is $\dfrac{R}{\sqrt2}$, where $R$ is the radius of earth. The surface of tunnel is assumed to be frictionless. A small particle of mass $m_0$ is released from point A. The collision between the particle and the tunnel is assumed to be perfectly inelastic. Take mass of earth is $M$. Radius of earth $R=6400$ km and $g=10$ m/s$^2$. Now answer the following Questions.
\qfig{p1_phy_s3_para5_6_1}
\qn{5} The time taken by the particle to move in the tunnel from A to O (in sec).
\qn{6} Work done by gravitational field of earth on the particle to reach from B to O is $x\dfrac{Gm_0M}{R}$. Find $x$.
\sect{SECTION-II(ii)}
\qn{7} A small elastic ball moving horizontally with 40 m/s collides elastically with the smooth stationary wedge of finite mass at point P which is placed on smooth horizontal surface as shown. After that it again strikes on the wedge at the same point P after $t_0$ seconds. Find the value of $t_0$ (in sec). (Take $g=10$ ms$^{-2}$)
\qfig{p1_phy_s4_q07_1}
\qn{8} Two identical potentiometer $P_1$ and $P_2$ of equal length $\ell$, resistance $R_1$, $R_2$ are connected with a battery of emf $\varepsilon_0$ and internal resistance $1\Omega$, through two switches $S_1$ and $S_2$. A battery of emf $\varepsilon$ is balanced on these potentiometer wires one by one. If resistance of wire $R_1$ is $3.68\ \Omega$ and balancing length is $\dfrac{\ell}{2}$ on it, when $S_1$ is closed and $S_2$ is open. On closing $S_2$ and opening $S_1$ the balancing length on $P_2$ is found to be $\dfrac{2\ell}{3}$, then the resistance of wire $R_2$ is $\alpha\Omega$. Fill the value of $16\alpha$.
\qfig{p1_phy_s4_q08_1}
\qn{9} A pendulum of mass $m$ connected with string of length $\ell=20\sqrt3$ cm. It is suspended above the centre of a sphere of radius 20 cm and mass $m$, at a distance 40 cm above centre of sphere. After collision pendulum stops. All surfaces are smooth. If coefficient of restitution is $e$. Write value of $100e$.
\qfig{p1_phy_s4_q09_1}
\partsec{2}{CHEMISTRY}
\sect{SECTION-I(i)}
\qn{1} (I) $\mathrm{Co(CO)_6}$\\
(II) $\mathrm{Co(CO)_5NH_3}$\\
(III) $\mathrm{Co(CO)_4(NH_3)_2}$\\
(IV) $\mathrm{Co(CO)_3(NH_3)_3}$\\
(V) $\mathrm{Co(CO)_2(NH_3)_4}$\\
(VI) $\mathrm{Co(CO)(NH_3)_5}$\\
Which statement is correct for these compounds?
\begin{qoptions}
\opt{A}{Order of Co--C bond length in complexes is I $>$ II $>$ III $>$ IV $>$ V $>$ VI}
\opt{B}{Order of Co--C bond length in complexes is VI $>$ V $>$ IV $>$ III $>$ II $>$ I}
\opt{C}{Order of Co--C bond strength in complexes is I $=$ II $=$ III $=$ IV $=$ V $=$ VI}
\opt{D}{Order of C--O bond length in complexes is I $=$ II $=$ III $=$ IV $=$ V $=$ VI}
\end{qoptions}
\qn{2} Which of the following reactant will not favour bimolecular nucleophilic substitution reaction?
\begin{qoptions}
\opt{A}{\qfigi{p1_chem_s1_q02_A}}
\opt{B}{\qfigi{p1_chem_s1_q02_B}}
\opt{C}{\qfigi{p1_chem_s1_q02_C}}
\opt{D}{All of above}
\end{qoptions}
\qn{3} \qfigi{p1_chem_s1_q03_1}
\begin{qoptions}
\opt{A}{\qfigi{p1_chem_s1_q03_A}}
\opt{B}{\qfigi{p1_chem_s1_q03_B}}
\opt{C}{\qfigi{p1_chem_s1_q03_C}}
\opt{D}{\qfigi{p1_chem_s1_q03_D}}
\end{qoptions}
\qn{4} 2 litres of a hard water sample having hardness due to $\mathrm{CaSO_4}$ is passed through anion exchange resin $[\mathrm{R(OH)_2}]$ and water coming out of resin required 5 ml 0.001 M $\mathrm{H_2SO_4}$ for neutralisation. Calculate hardness of this sample in terms of $\mathrm{CaCO_3}$. [Atomic mass : Ca $=40$, S $=32$]
\begin{qoptions}
\opt{A}{$62.5$}
\opt{B}{$125$}
\opt{C}{$250$}
\opt{D}{$500$}
\end{qoptions}
\sect{SECTION-I(ii)}
\qn{5} Which of the following statement(s) is(are) false?
\begin{qoptions}
\opt{A}{In $\mathrm{XeOF_4}$, the central atom Xe doesn't lie on the plane of the square}
\opt{B}{In $\mathrm{XeF_6}$ $d_{x^2-y^2}$, $d_{z^2}$ and $d_{xy}$ orbitals are used in bonding}
\opt{C}{Maleic acid is more acidic as compared to fumaric acid}
\opt{D}{The dipole moment of $\mathrm{CH_3F}$ is greater than $\mathrm{CH_3Cl}$}
\end{qoptions}
\qn{6} Choose the correct statement(s).
\begin{qoptions}
\opt{A}{$\mathrm{NF_3}$ undergoes hydrolysis through $\mathrm{S_N1}$ mechanism under high temperature}
\opt{B}{Rate of hydrolysis $\mathrm{SiCl_4}>\mathrm{SiF_4}$}
\opt{C}{Upon hydrolysis of $\mathrm{BeCl_2}$ through $\mathrm{S_N2}$ mechanism produces a final product which dissolves in exess NaOH}
\opt{D}{$\mathrm{SbCl_3}$ produces white turbidity of $\mathrm{Sb(OH)_3}$, upon hydrolysis.}
\end{qoptions}
\qn{7} In which of the following species, when $1\,e^-$ is added it must entered in gerade molecular orbital and after adding that $1\,e^-$, total electrons in ungerade molecular orbitals is 8
\begin{qoptions}
\opt{A}{$\mathrm{N_2}$}
\opt{B}{$\mathrm{O_2}$}
\opt{C}{$\mathrm{F_2}$}
\opt{D}{$\mathrm{C_2}$}
\end{qoptions}
\qn{8} $\mathrm{CH_3{-}CH{=}CH_2}\xrightarrow{\mathrm{Cl_2/aq.NaBr}}$\\
Product can be :
\begin{qoptions}
\opt{A}{$\mathrm{CH_3{-}CH(OH){-}CH_2Cl}$}
\opt{B}{$\mathrm{CH_3{-}CH(Br){-}CH_2Cl}$}
\opt{C}{$\mathrm{CH_3{-}CH(Cl){-}CH_2Br}$}
\opt{D}{$\mathrm{CH_3{-}CH(Cl){-}CH_2Cl}$}
\end{qoptions}
\qn{9} Consider the given statements and select the correct option(s) for HCP unit cell :
\begin{qoptions}
\opt{A}{Packing efficiency is more when all tetrahedral voids are filled with spheres of maximum possible size as compared with packing efficiency when all octahedral voids are filled in similar way.}
\opt{B}{Interplanar spacing between closely packed planes is $\left(2r\sqrt{\dfrac23}\right)$, where $r$ = radius of atom.}
\opt{C}{The hexagonal base of HCP unit cell can be divided into three identical rhombuses}
\opt{D}{There are a total of 6 octahedral voids effectively per unit cell.}
\end{qoptions}
\qn{10} Which of the following statement(s) is/are correct ? ($R$ = atomic radius)
\begin{qoptions}
\opt{A}{In CCP, distance between two nearest octahedral voids is $2R$}
\opt{B}{In CCP, distance between nearest tethrahedral and octahedral voids is $2\sqrt3R$}
\opt{C}{In CCP, distance between two nearest voids is $\sqrt2R$}
\opt{D}{In CCP, distance between next nearest tetrahedral void is $\sqrt5R$}
\end{qoptions}
\sect{SECTION-II(i)}
\paraq{1}{2}
Ligands can be classified according to their electron donating and accepting tendency as classical and non classical ligands.
\qn{1} Among the following ligands, how many are non classical ligands and donates electron through $\sigma$ bond and accepts in gerade molecular orbital?\\
$\mathrm{CO}$, $\mathrm{NO}$, $\mathrm{C_2H_4}$, $\mathrm{NH_3}$, $\mathrm{NH_2^-}$, $\mathrm{Br^-}$, $\mathrm{O^{2-}}$, $\mathrm{N_2}$, $\mathrm{PCl_3}$, $\mathrm{C_5H_5^-}$, $\mathrm{C_6H_6}$, $\mathrm{NO^{\oplus}}$
\qn{2} How many of the following statements is(are) correct?\\
(i) Stability order : $\mathrm{[Mn(CO)_6]<[Mn(CO)_6]^+}$\\
(ii) Stability order : $\mathrm{[Co(CO)_4]<[Co(CO)_4]^-}$\\
(iii) The formula for the IUPAC name dicarbonyldinitrosyl Iron(0) is $\mathrm{[Fe(CO)_2(NO)_2]}$\\
(iv) EAN value for Iron in diIron nonacarbonyl is 35\\
(v) Total number of coordination position isomers for $\mathrm{[Co(NH_3)_4(NH_2)OHCo(H_2O)_4]^{4+}}$ is 5
\paraq{3}{4}
Addition of HX on alkene proceed through the formation of carbocation. This reaction is also known as Markownikoff reaction. According to Markownikoff rule addition of electrophile occurs on that carbon of alkene which has more number of H atom.
\qfig{p1_chem_s3_para3_4_1}
\qn{3} How many total possible products (including stereoisomer) can be formed in given reaction?\\
$\mathrm{CH_3{-}CH{=}CH{-}C_2H_5}\xrightarrow{\mathrm{HCl}}$
\qn{4} Number of alkene which are more reactive then propene with HCl / $\mathrm{CCl_4}$ ?\\
(1) $\mathrm{CH_2{=}CH_2}$ \quad (2) $\mathrm{Ph{-}CH{=}CH_2}$\\
(3) $\mathrm{Cl_3C{-}CH{=}CH_2}$ \quad (4) $\mathrm{CH_2{=}CH{-}NO_2}$\\
(5) $\mathrm{CH_2{=}CH{-}OCH_3}$ \quad (6) $\mathrm{CH_2{=}CH{-}NMe_2}$
\paraq{5}{6}
Vapour pressure of solution containing 12g of a non-volatile, non electrolyte solute in 180 g water is 20.0 Torr. If 1 mole water is further added then vapour pressure of solution increases by 1 Torr. (Consider moles of solute to be much less than moles of water in solution)
\qn{5} What is the molar mass of solute (in g mol$^{-1}$)
\qn{6} What is the vapour pressure of pure water is (in Torr)
\sect{SECTION-II(ii)}
\qn{7} How many of the following ligands form atleast 1 five membered ring having atleast 1 nitrogen and oxygen as donor atoms :\\
Oxalato (ox)$^{-2}$, Acetylacetonato (acac)$^{-1}$, Ethylene diamine (en), Propylene diamine (pn), Butylene diamine (bn), Imidodiacetate (imda$^{2-}$), NTA$^{3-}$, EDTA$^{3-}$, DMG$^{-1}$, gly$^{-1}$
\qn{8} Total number of geometrical isomers are possible for following compound :
\qfig{p1_chem_s4_q08_1}
\qn{9} In cubical close packing (ccp) arrangement of solid what will be total effective number of voids present in the given plane in the given diagram. If your answer is $x$ then write $100x$.
\qfig{p1_chem_s4_q09_1}
\partsec{3}{MATHEMATICS}
\sect{SECTION-I(i)}
\qn{1} Number of solutions of the equation $\left|\tan^{-1}|x|\right|=\sqrt{\left(x^2+\dfrac12\right)^2-2x^2}$ is
\begin{qoptions}
\opt{A}{$1$}
\opt{B}{$2$}
\opt{C}{$3$}
\opt{D}{$4$}
\end{qoptions}
\qn{2} The total number of solutions of the equation $[x]^2=x+2\{x\}$, (where $[.]$, $\{.\}$ denotes greatest integer function and fractional part function respectively)
\begin{qoptions}
\opt{A}{$2$}
\opt{B}{$4$}
\opt{C}{$6$}
\opt{D}{$5$}
\end{qoptions}
\qn{3} The value of {\scriptsize$\displaystyle\lim_{x\to0}\dfrac{(1+3x+2x^2)^{1/x}-(1+3x-2x^2)^{1/x}}{x}$} is
\begin{qoptions}
\opt{A}{$-e^3$}
\opt{B}{$3e^2$}
\opt{C}{$\dfrac13e^{-2}$}
\opt{D}{$4e^3$}
\end{qoptions}
\qn{4} Positive integers $p$, $q$ and $r$ satisfy the condition that the quadratic equation $px^2-qx+r=0$ has two distinct real roots that are strictly between 0 and 1. Then the minimum value of $p$, is
\begin{qoptions}
\opt{A}{$2$}
\opt{B}{$1$}
\opt{C}{$5$}
\opt{D}{$4$}
\end{qoptions}
\sect{SECTION-I(ii)}
\qn{5} The value of $\displaystyle\sum_{r=2}^{100}\dfrac{2^r(2-r)}{r(r+1)(r+2)}$ is equal to
\begin{qoptions}
\opt{A}{$\dfrac23-\dfrac{2^{100}}{(51)(101)}$}
\opt{B}{$\dfrac23-\dfrac{2^{101}}{(101)(102)}$}
\opt{C}{$\dfrac12-\dfrac{2^{100}}{100.101}$}
\opt{D}{$\dfrac12-\dfrac{2^{100}}{(101)(51)}$}
\end{qoptions}
\qn{6} Consider the function {\scriptsize$f(x)=\displaystyle\lim_{n\to\infty}\dfrac{(1+\cos x)^n+5\log_ex}{2+(1+\cos x)^n}$}, then
\begin{qoptions}
\opt{A}{$f(x)$ is continuous at positive odd integral multiple of $\pi$}
\opt{B}{$f(x)$ is discontinuous at positive even integral multiple of $\pi$}
\opt{C}{$f(x)$ is discontinuous at positive odd integral multiple of $\dfrac\pi2$}
\opt{D}{$f(x)$ is continuous at positive even integral multiples of $\dfrac\pi2$}
\end{qoptions}
\qn{7} If {\footnotesize$P_n=\displaystyle\lim_{x\to0^+}$}\par\noindent\hfill{\scriptsize$\left(\log(\tan x)+\log(\tan2x)+\log(\tan3x)+\ldots+\log(\tan nx)-\log x^n\right)$}\hfill\null\par\par\noindent and $P_{256}-P_{255}=kP_m$, (where $k,m,n\in N$), then-
\begin{qoptions}
\opt{A}{$k=8$}
\opt{B}{$k=6$}
\opt{C}{$m=2$}
\opt{D}{$m=4$}
\end{qoptions}
\qn{8} Let $f(x)=\begin{cases}-2 & \text{for}-3\leqslant x\leqslant0\\ x-2 & \text{for}\ 0<x\leqslant3\end{cases}$, $g(x)=[f(|x|)]+[|f(x)|]$ where $[.]$ denotes the greatest integer function. Which of the following statement(s) is/are true ?
\begin{qoptions}
\opt{A}{$g(1.7)=-1$}
\opt{B}{$g(x)$ is not differentiable at any point}
\opt{C}{$g(x)$ is continuous at all non integral points}
\opt{D}{$g(x)$ is differentiable at some points}
\end{qoptions}
\qn{9} Let $f:\left[-\dfrac13,3\right]\to R$ and $g:\left[-\dfrac13,3\right]\to R$ be function defined by $f(x)=[x^2-4]$ and $g(x)=|x-2|\,f(x)+|3x-5|\,f(x)$, where $[y]$ denotes the greatest integer less than or equal to $y$ for $y\in R$. Then
\begin{qoptions}
\opt{A}{$f$ is discontinuous exactly at eight points in $\left[-\dfrac13,3\right]$}
\opt{B}{$f$ is discontinuous exactly at nine points in $\left[-\dfrac13,3\right]$}
\opt{C}{$g$ is NOT differentiable exactly at ten points in $\left(-\dfrac13,3\right)$}
\opt{D}{$g$ is NOT differentiable exactly at nine points in $\left(-\dfrac13,3\right)$}
\end{qoptions}
\qn{10} If $\alpha$ ($\alpha>0$) is a root of $x^{11}-x^{10}+x^8-x^7+x^5-x^4+x^2-x-20=0$, then $\alpha^{12}$ is greater than
\begin{qoptions}
\opt{A}{$59$}
\opt{B}{$61$}
\opt{C}{$58$}
\opt{D}{$57$}
\end{qoptions}
\sect{SECTION-II(i)}
\paraq{1}{2}
Let $g:R\to\left\{\dfrac12\right\}$ be a function given by $g(x)=x^{10}(f'(t)-2)+x^9(f''(t))+4xf(0)+\dfrac12$ where $t\in R$.
\qn{1} Number of solutions of $|f(x)|=\cos x$, is -
\qn{2} Number of points of non differentiability of function $h(x)=\left|f(x)(e^{f(x)}-1)(f(x)-1)(f(x)-2)\right|$, is
\paraq{3}{4}
Consider the function $f:A\to[0,5]$, $f(x)=\sqrt{9-(x-1)(x-3)(x-5)(x-7)}$.\\
(Where set $A$ represents all values of $x$ for which $f(x)$ is defined.)
\qn{3} In the domain of the function, number of possible integral values of $x$ is equal to
\qn{4} Among the following, the number of correct statements is\\
(i) function is many one \quad (ii) function is one-one\\
(iii) function is onto \quad (iv) function is into
\paraq{5}{6}
Let $g(x)=6x^2-18x+8$;\\
$f_1(x)=|g(x)|$, $f_2(x)=|f_1(x)-P_1|$, $f_3(x)=|f_2(x)-P_2|$, $f_4(x)=|f_3(x)-P_3|$.
\qn{5} If $P_1=7$, then the largest integral value of $P_2$ such that $f_3(x)$ has exactly 10 points of non-differentiability, is
\qn{6} If $P_1,P_2,P_3$ are consecutive primes such that $5<P_1<P_2<P_3$ and total number of non-differentiable points of $f_1,f_2,f_3$ are $a_1,a_2,a_3$ respectively, then the number of solution of equation $H(x)=\dfrac{(x-a_1)(x-a_3)}{(x-a_2)}=1$, is
\sect{SECTION-II(ii)}
\qn{7} Let $A=\{1,2,3,4\}$, $B=\{1,2,3,4\}$. The total number of functions $f:A\to B$ such that $f(x)\neq x\ \forall\ x\in A$ and range of $f(x)$ contains exactly three elements, is
\qn{8} Let $T_n$ be the $n^{\text{th}}$ term of a series for $n\geq1$ and $T_1=1$, $T_{n+1}=2T_n+4^n$, $n\geq1$. If $S_n=$ Sum of first $n$ terms of the series and $S_n=\dfrac{a.4^n+b.2^n+1}{3}$ where $a,b\in I$, then the value of $10\left|\dfrac ba\right|$ is
\qn{9} If $\displaystyle\lim_{x\to-\infty}\left(\sqrt{x^2-\dfrac{4x}{3}}+x+2\right)$ is equal to $K$ then value of $6K$ is-
\papersec{2}
\partsec{1}{PHYSICS}
\sect{SECTION-I(i)}
\qn{1} In the given circuit switch is closed at $t=0$ and capacitor is initially uncharged. Graph between $V_B-V_A$ and time is shown for different values of variable resistance $R$. Which are correct options.
\qfig{p2_phy_s1_q01_1}
\qfig{p2_phy_s1_q01_2}
\begin{qoptions}
\opt{A}{$|V_1|>|V_2|$}
\opt{B}{$|V_1|<|V_2|$}
\opt{C}{$R_1>R_2$}
\opt{D}{$R_1<R_2$}
\end{qoptions}
\qn{2} A galvanometer of resistance $10\Omega$ is shunted with a series combination of resistance $3\Omega$ and cell of emf $4$V as shown. Current for full scale deflection in galvanometer is $2.5$ A. Choose the correct option(s).
\qfig{p2_phy_s1_q02_1}
\begin{qoptions}
\opt{A}{Current through galvanometer is $2$A.}
\opt{B}{Current through galvanometer is $1$A.}
\opt{C}{Maximum value of $I$ can be $12.17$A.}
\opt{D}{Maximum value of $I$ can be $15.37$A.}
\end{qoptions}
\qn{3} A spherical conducting shell having total charge $10$ C is fixed inside a uniform electric field of $40$ N/C. Net force on $4$C charge placed inside shell at distance $R/2$ from centre is $-60\,\mathrm N\,\hat j$. Choose incorrect option(s).
\qfig{p2_phy_s1_q03_1}
\begin{qoptions}
\opt{A}{Force on $4$C due to external electric field is zero.}
\opt{B}{Force on sphere due to external electric field is $560$ N.}
\opt{C}{Force on $4$C charge particle due to charge on outer surface of spherical shell is $-160\,\mathrm N\,\hat j$.}
\opt{D}{Force on $4$C charge due to spherical shell is $220\,\mathrm N\,\hat j$.}
\end{qoptions}
\qn{4} Two stars are undergoing circular motion about their common centre of mass. Motion observed by an observes which is very far and in the plane of motion is SHM. Velocity versus phase rotated by both A and B is as shown. $R_1=8\times10^7$ km. (Take $G=\dfrac{20}{3}\times10^{-11}$ SI)
\qfig{p2_phy_s1_q04_1}
\qfig{p2_phy_s1_q04_2}
\begin{qoptions}
\opt{A}{Distance between both stars is $10^{11}$ m.}
\opt{B}{Distance between both stars is $10^8$ m.}
\opt{C}{Total mass of both stars is $\dfrac38\times10^{25}$ kg.}
\opt{D}{Total mass of both stars is $\dfrac78\times10^{25}$ kg.}
\end{qoptions}
\qn{5} A parallel plate capacitor is partially filled with dielectric. Its potential with respect to position graph is as shown. Area of plates is $2\,\mathrm m^2$.
\qfig{p2_phy_s1_q05_1}
\begin{qoptions}
\opt{A}{Dielectric constant is $2$.}
\opt{B}{Dielectric constant is $3$.}
\opt{C}{Capacitance is $\dfrac{2000\varepsilon_0}{9}$.}
\opt{D}{Capacitance is $\dfrac{4000\varepsilon_0}{9}$.}
\end{qoptions}
\qn{6} Two charged particles M and N enter a space of a uniform magnetic field with velocities perpendicular to the magnetic field. The paths are as shown in Figure below. Momentum of the charged particles are $p_1$ and $p_2$ respectively, then :-
\qfig{p2_phy_s1_q06_1}
\begin{qoptions}
\opt{A}{the charge of M is greater than that of N if $p_1=p_2$}
\opt{B}{the momentum of N is greater than that of M if $q_M=q_N$.}
\opt{C}{the specific charge of M is greater than that of N, if $v$ is same}
\opt{D}{the speed of M is less than that of N if specific charge is same}
\end{qoptions}
\sect{SECTION-I(ii)}
\paraq{7}{8}
An infinite line charge having initial charge density $\lambda_0$ is placed along the axis of a infinite conducting cylinder made of material having resistivity $\rho$ and permittivity $\varepsilon$.
\qn{7} What is the line charge density on wire at time $t$?
\begin{qoptions}
\opt{A}{$\dfrac{\lambda_0}{2}e^{\frac{-2t}{\varepsilon\rho}}$}
\opt{B}{$\lambda_0e^{-\frac{t}{\varepsilon\rho}}$}
\opt{C}{$\lambda_0-\dfrac{\varepsilon}{\rho}t$}
\opt{D}{$\lambda_0e^{-\frac{2t}{\varepsilon\rho}}$}
\end{qoptions}
\qn{8} What will be current density at radial distance $r$?
\begin{qoptions}
\opt{A}{$\dfrac{\lambda\rho}{2\pi\varepsilon r}$}
\opt{B}{$\dfrac{\lambda r}{2\pi\varepsilon\rho}$}
\opt{C}{$\dfrac{\lambda}{2\pi\varepsilon\rho r}$}
\opt{D}{None}
\end{qoptions}
\paraq{9}{10}
A system consists of two small spheres of mass 1kg and 2 kg interconnected by a weightless spring. At $t=0$, the spheres are set in motion with initial velocities as shown in figure. Position vector of centre of mass at $t=0$ is $20\hat j$. Initially spring is in its natural length. System moves in earth's gravitational field.
\qfig{p2_phy_s2_para9_10_1}
\qn{9} The momentum of the system at $t=2$ sec in kg m/s is ($g=10$ m/s$^2$)
\begin{qoptions}
\opt{A}{$\left(12\hat i-32\hat j\right)$}
\opt{B}{$\left(16\hat i-64\hat j\right)$}
\opt{C}{$\left(18\hat i-64\hat j\right)$}
\opt{D}{$\left(12\hat i-18\hat j\right)$}
\end{qoptions}
\qn{10} The position vector of centre of mass of the system at $t=1$s is :-
\begin{qoptions}
\opt{A}{$4\hat i+\dfrac{24}{25}\hat j$}
\opt{B}{$6\hat i+\dfrac{39}{2}\hat j$}
\opt{C}{$4\hat i+\dfrac{29}{3}\hat j$}
\opt{D}{$6\hat i+\dfrac{41}{3}\hat j$}
\end{qoptions}
\sect{SECTION-II(i)}
\paraq{1}{2}
A charge particle $q$ is held between two conducting shells as shown. Both of conducting shell are neutral. (Take $q=2\,\mu$C, $R=20$ m)
\qfig{p2_phy_s3_para1_2_1}
\qn{1} Sum of potential (in V) of both spheres will be :-
\qn{2} If outer surface is grounded by connecting with a resistance $R$, then heat loss (in mJ) in resistance $R$ is
\paraq{3}{4}
Consider a simple RC circuit as shown in figure and capacitor is initially at steady state. If by some means capacitance is suddenly reduced to one third of its initial capacitance. Initially capacitance is $2\mu$F.
\qfig{p2_phy_s3_para3_4_1}
\qn{3} Find current (A) in circuit just after capacitance is reduced.
\qn{4} Find magnitude of work done by battery till steady state is achieved again in $\mu$J.
\paraq{5}{6}
The melting point of wax is measured by using a thermometer which is not exact, and it indicates a temperature of $50^\circ$C as shown in figure-1.
\qfig{p2_phy_s3_para5_6_1}
A wheat stone bridge is used to measure this melting point exactly. The deviation from $50^\circ$C is being measured by observing the variation of the resistance of a coil embedded in wax.\\
The wheat stone bridge is setup so that the resistance of each arm is $50\ \Omega$ exactly at $50^\circ$C and there is no current in the galvanometer as shown in figure-II. The resistance of the galvanometer is $50\ \Omega$.
\qfig{p2_phy_s3_para5_6_2}
One arm of the bridge containing the coil of resistance $50\ \Omega$ is now immersed in a wax-box. The wax is just melted as shown in the figure-III. It is now observed that the current through the galvanometer is $0.1\ \mu$A. Given, temperature coefficient of resistance is $10^{-3}$ per $^\circ$C. Resistance other than inside wax-box are assumed to be constant and are exactly $50\Omega$ while the resistance of wax box is very close to $50\Omega$.
\qfig{p2_phy_s3_para5_6_3}
\qn{5} The variation in the resistance (in m$\Omega$) inside the wax-box from its initial value of $50\Omega$ is :-
\qn{6} The deviation (in $^\circ$C) in the melting point from $50^\circ$C is
\sect{SECTION-II(ii)}
\qn{7} A trolley is moving with constant velocity of 1 m/s on a straight track due to some external force of 20N. Sand is continuously falling on it with negligible velocity and moves with trolley. What is magnitude of rate of work done by friction between sand and trolley?
\qn{8} A semi-cylinder is cut out of a block and the block is kept on the horizontal surface. A small body of mass $m$ is released at the top of the cylinder shaped hole. (Friction is negligible).\\
If the block is able to move on the horizontal surface without friction then the force exerted on the block by the small body at the lowest point is $7$ mg/2. Find the ratio of $M/m$ ?
\qfig{p2_phy_s4_q08_1}
\qn{9} In the show circuit all cells are ideal. If current in branch BE is $x$ ampere and current in branch AE is $y$ ampere then value of $x+y$ (in ampere) is
\qfig{p2_phy_s4_q09_1}
\partsec{2}{CHEMISTRY}
\sect{SECTION-I(i)}
\qn{1} Reactants which gives $\mathrm{S_N2}$ reaction with KCN in DMSO solvent?
\begin{qoptions}
\opt{A}{$\mathrm{Me_3C{-}CH_2{-}Cl}$}
\opt{B}{\qfigi{p2_chem_s1_q01_B}}
\opt{C}{\qfigi{p2_chem_s1_q01_C}}
\opt{D}{\qfigi{p2_chem_s1_q01_D}}
\end{qoptions}
\qn{2} \qfigi{p2_chem_s1_q01_4}\\
(X) may be :
\begin{qoptions}
\opt{A}{\qfigi{p2_chem_s1_q02_A}}
\opt{B}{\qfigi{p2_chem_s1_q02_B}}
\opt{C}{\qfigi{p2_chem_s1_q02_C}}
\opt{D}{\qfigi{p2_chem_s1_q02_D}}
\end{qoptions}
\qn{3} Structure of $\mathrm{FXeN(SO_2F)_2}$ (pale yellow solid) and $\mathrm{FXeOSO_2F}$ are given below\\
(It is observed that geometry of nitrogen is planar with its surrounding atom in structure-I)
\qfig{p2_chem_s1_q03_1}
Select the \textbf{CORRECT} statements?
\begin{qoptions}
\opt{A}{Maximum seven atoms may lie in one plane in structure (I)}
\opt{B}{10 atoms are $sp^3$ hybridized in structure (I)}
\opt{C}{Nine d-orbitals are involved in bonding in structure (I)}
\opt{D}{S--O single bond length of $\mathrm{FXeOSO_2F}$ is lesser than S--N bond length in $\mathrm{FXeN(SO_2F)_2}$}
\end{qoptions}
\qn{4} Which of the following set of order is \textbf{INCORRECT} for their indicated properties?
\begin{qoptions}
\opt{A}{Covalent character : $\mathrm{CuCl<AgCl<AuCl}$}
\opt{B}{Stability : $\mathrm{PbF_4>PbCl_4>PbBr_4>PbI_4}$}
\opt{C}{Acidic nature order : $\mathrm{BeO<MgO<CaO<SrO}$}
\opt{D}{Water solubility : {\small$\mathrm{LiNO_3<NaNO_3<KNO_3<RbNO_3}$}}
\end{qoptions}
\qn{5} Choose incorrect statement(s) is/are :
\begin{qoptions}
\opt{A}{In presence of non-stoichiometric defects, ionic crystals lost their electrical neutrality.}
\opt{B}{$e^-$ holes or $e^-$ vacancy move towards positively charged plate, under influence of electric field in p-type semiconductor.}
\opt{C}{$a=b\neq c$ represents both tetragonal \& trigonal crystal system.}
\opt{D}{HCP and CCP both have 74\% packing efficiency.}
\end{qoptions}
\qn{6} A solution prepared by mixing 1 mol benzene ($P^\circ_{\text{benzene}}=42$ mm) and 2 mol toluene ($P^\circ_{\text{toluene}}=36$ mm) will have :
\begin{qoptions}
\opt{A}{Total vapour pressure of 38 mm.}
\opt{B}{Mole fraction of vapour of benzene above liquid mixture is $7/19$.}
\opt{C}{Benzene is more volatile than toluene}
\opt{D}{Toluene has higher normal boiling point than benzene}
\end{qoptions}
\sect{SECTION-I(ii)}
\paraq{7}{8}
When a substance undergoes nucleophilic substitution reaction and the nucleophile is the solvent itself, then the reaction is known as solvolysis. In solvolysis reaction, if the solvent used is water then the reaction is called as hydrolysis.
\qn{7} Select the INCORRECT match regarding hydrolysis.
\begin{qoptions}
\opt{A}{$\mathrm{P_4O_{10}}$ : During partial hydrolysis oxy acid form having basicity $=6$}
\opt{B}{$\mathrm{SiF_4}$ \& $\mathrm{BF_3}$ : Partial hydrolysis}
\opt{C}{$\mathrm{CF_4}$ \& $\mathrm{NF_3}$ : Do not hydrolysed}
\opt{D}{$\mathrm{Mg_2C_3}$ \& $\mathrm{Al_4C_3}$ : Same gas release on hydrolysis}
\end{qoptions}
\qn{8} Which of the following show disproportionation redox reaction in hydrolysis.
\begin{qoptions}
\opt{A}{$\mathrm{XeF_4}$ \& $\mathrm{XeF_6}$}
\opt{B}{$\mathrm{N_2O_4}$ \& $\mathrm{PCl_3}$}
\opt{C}{$\mathrm{ClF_3}$ \& $\mathrm{BrF_5}$}
\opt{D}{$\mathrm{XeF_4}$ and $\mathrm{NO_2}$}
\end{qoptions}
\paraq{9}{10}
Law of equivalence states that whenever two substances react, the equivalents of the two reactants will be equal and the equivalents of the product(s) formed will also be equal. This law finds great application in acid-base as well as in redox reactions.
\qn{9} 50 ml of an aqueous solution containing $a$ M $\mathrm{H_2SO_4}$ and $b$ M $\mathrm{H_2C_2O_4}$ required 25 ml of 1 M $\mathrm{Ca(OH)_2}$ for complete neutralisation. The same solution separately required 10 ml of 0.6 M acidified $\mathrm{KMnO_4}$ solution for complete titration. Select the correct combination of $a$ and $b$.
\begin{qoptions}
\opt{A}{$a=0.1$, $b=0.4$}
\opt{B}{$a=0.3$, $b=0.2$}
\opt{C}{$a=0.2$, $b=0.2$}
\opt{D}{$a=0.2$, $b=0.3$}
\end{qoptions}
\qn{10} Iron pyrites ($\mathrm{FeS_2}$) upon roasting produces $\mathrm{SO_2}$ and $\mathrm{Fe_2O_3}$. The amount of $\mathrm{FeS_2}$ ($X$ grams) required such that the $\mathrm{SO_2}$ produced completely decolorises 1 L of 0.1 M $\mathrm{KMnO_4}$ solution. Unbalanced equations are :\\
$\mathrm{FeS_2+O_2\to Fe_2O_3+SO_2}$\\
$\mathrm{KMnO_4+SO_2\to MnSO_4+K_2SO_4+H_2SO_4}$\par
Select the incorrect option from the following :\par
{[}Atomic mass : Fe $=56$, S $=32$, Mn $=55$, K $=39${]}
\begin{qoptions}
\opt{A}{Equivalent weight of $\mathrm{SO_2}$ in the $2^{\text{nd}}$ reaction is 32 g/eq.}
\opt{B}{The number of electrons gained per formula unit of $\mathrm{KMnO_4}$ is 5}
\opt{C}{Number of electron lost per formula unit of $\mathrm{FeS_2}$ is 11}
\opt{D}{$X=30$ g}
\end{qoptions}
\sect{SECTION-II(i)}
\paraq{1}{2}
Acid catalysed conversion of 1,2-diol or vicinal diol, into carbonyl compound known as pinacol-pinacolone rearrangement.
\qfig{p2_chem_s3_para1_2_1}
Generally more electron donating group migrate during the mechanism.
\qn{1} Calculate the degree of unsaturation in the final product of given reaction?
\qfig{p2_chem_s3_q01_1}
\qn{2} Number of lone pair present in the major product of the given reaction
\qfig{p2_chem_s3_q02_1}
\paraq{3}{4}
The mechanism of acid catalysed dehydration is given below
\qfig{p2_chem_s3_para3_4_1}
This reaction take place in presence of conc. $\mathrm{H_2SO_4}$. In this reaction carbocation is formed and rearrangement may be possible.
\qn{3} Calculate the number of total alkene formed by dehydration of 2-butanol?
\qn{4} How many 1,2-shifts of carbocation intermediate are involved during the course of following reaction?
\qfig{p2_chem_s3_q04_1}
\paraq{5}{6}
Compounds that have the same chemical formula but different structural arrangements are called isomers. Because of the complicated formulae of many coordination compound, the variety of bond types and the number of shapes possible, many different types of isomerism occur. Werner's classification into polymerization, ionization, hydrate linkage, coordination, coordination position, and geometric and optical isomerism is still generally accepted.
\qn{5} Find total number of structural isomers possible for complex. [Consider monohydrate \& dihydrate isomers also]\\
$\mathrm{[Co(NH_3)_3(H_2O)_2Cl](NO_2)_2}$
\qn{6} How many total coordination isomers possible for complex $\mathrm{[Co(NH_3)_3(H_2O)_3]\,[Cr(CN)_6]}$.
\sect{SECTION-II(ii)}
\qn{7} Total product formed by mono chlorination of methyl cyclobutane?
\qfig{p2_chem_s4_q07_1}
\qn{8} Maximum number of atoms present in same plane for given molecules $\mathrm{N(SiH_3)_3}$ and $\mathrm{Al_2(CH_3)_6}$ are $x$ and $y$ respectively then calculate $(x+y)$.
\qn{9} A solution of a non-volatile solute in water freezes at 272.7 K. The vapour pressure of pure water at 298 K is 23.51 mm Hg and $K_f$ for water is 1.86 K kg mol$^{-1}$. Calculate the vapour pressure of this solution (in Torr) at 298 K.\\
{[}Freezing point of water $=273$ K{]}\\
{[}Write your answer excluding decimal places{]}
\partsec{3}{MATHEMATICS}
\sect{SECTION-I(i)}
\qn{1} If $a,b,c$ are positive real numbers such that
{\small\[\dfrac{1}{1+a^4}+\dfrac{1}{1+b^4}+\dfrac{1}{1+c^4}=\dfrac{\sqrt{1+a^4}+\sqrt{1+b^4}+\sqrt{1+c^4}}{\sqrt{(1+a^4)(1+b^4)(1+c^4)}},\]}
then which of the following statement(s) is/are true ?
\begin{qoptions}
\opt{A}{$\dfrac ab+\dfrac bc+\dfrac ca\leqslant3$}
\opt{B}{$\dfrac ab+\dfrac bc+\dfrac ca\geqslant3$}
\opt{C}{$(a+b+c)\left(\dfrac1a+\dfrac1b+\dfrac1c\right)\leqslant9$}
\opt{D}{$(a+b+c)\left(\dfrac1a+\dfrac1b+\dfrac1c\right)\geqslant9$}
\end{qoptions}
\qn{2} Let $f(x)$ be a twice differentiable function satisfying $xf(x^2)+(x^2+1)f(x)=x$, $\forall\ x\in R$ then the value of $f''(1)$ is
\begin{qoptions}
\opt{A}{$\dfrac13$}
\opt{B}{$-\dfrac13$}
\opt{C}{$-\dfrac19$}
\opt{D}{$1$}
\end{qoptions}
\qn{3} Identify the CORRECT statements-
\begin{qoptions}
\opt{A}{Sum of two periodic functions may or maynot be a periodic function}
\opt{B}{A differentiable function in $R$ will always be continuous in $R$.}
\opt{C}{A continuous function whose domain is in a closed interval is always bounded}
\opt{D}{In $[a,b]$ for a function $f(x)$, if $f(a).f(b)<0$, then there exists atleast one root of $f(x)=0$ in $(a,b)$.}
\end{qoptions}
\qn{4} If $f(x)=\displaystyle\lim_{n\to\infty}\dfrac{2x^{2n}+x+2}{x^{2n}-x^2+3}$, then which of the following statement(s) holds good ?
\begin{qoptions}
\opt{A}{$f(x)$ is continuous at $x=1$}
\opt{B}{$f(x)$ is discontinuous at $x=1$}
\opt{C}{$\displaystyle\lim_{x\to2}f(x)$ exists finitely}
\opt{D}{$\displaystyle\lim_{x\to2}f(x)$ does not exist}
\end{qoptions}
\qn{5} Let $f:[-1,1]\to[-1,1]$, $f(x)=\sin(2\sin^{-1}x)$ \& let $g(x)=\dfrac{f(x)}{\sqrt{1-x^2}}$, then -
\begin{qoptions}
\opt{A}{$f(x)$ is one-one function}
\opt{B}{$f(x)$ is odd function}
\opt{C}{$g(x)$ is continuous in $(-1,1)$}
\opt{D}{$f(x)$ is continuous for all $x\in[-1,1]$}
\end{qoptions}
\qn{6} $\displaystyle\lim_{x\to\infty}\sin^{-1}\sin\left(\dfrac{x^2+2}{x^2+1}\right)$ is equal to (where $[.]$ and $\{.\}$ denote greatest integer and fractional part function respectively)
\begin{qoptions}
\opt{A}{$\displaystyle\lim_{x\to\infty}\left[\dfrac{x^2+2}{x^2+1}\right]$}
\opt{B}{$\displaystyle\lim_{x\to\infty}\dfrac{x^2+2}{x^2+1}$}
\opt{C}{$\displaystyle\lim_{x\to\infty}\left\{\dfrac{x^2+2}{x^2+1}\right\}$}
\opt{D}{$\displaystyle\lim_{x\to\infty}\left\{\dfrac{x^2}{x^2+1}\right\}$}
\end{qoptions}
\sect{SECTION-I(ii)}
\paraq{7}{8}
Consider $f(x)=\left[\ell n\left(\sin x+\sqrt{\sin^2x+1}\right)\right]+\left[\ell n\left(\sqrt{\sin^2x+1}-\sin x\right)\right]$\\
$g(x)=\left\{\dfrac x\pi\right\}-\left[\dfrac{\sin^22x+\cos^42x}{\cos^22x+\sin^42x}\right]$ and $h(x)=f(x)+g(x)$\\
(where $[.]$ and $\{.\}$ are greatest integer and fractional part function respectively)
\qn{7} Which of following is correct -
\begin{qoptions}
\opt{A}{Equation $f(x)=g(x)$ has infinite many solutions}
\opt{B}{Period of $f(x)$ is not defined}
\opt{C}{$\displaystyle\lim_{x\to0}f(x)$ does not exist}
\opt{D}{$\displaystyle\lim_{x\to0}f(x).g(x)$ does not exist}
\end{qoptions}
\qn{8} Number of points in $[-2\pi,2\pi]$ where $h(x)$ is not differentiable, is-
\begin{qoptions}
\opt{A}{$2$}
\opt{B}{$3$}
\opt{C}{$4$}
\opt{D}{$5$}
\end{qoptions}
\paraq{9}{10}
Let $a$ is the number of points of continuity of $f_1(x)=\begin{cases}x-1 & x\in Q\\ x^2-x-2 & x\notin Q\end{cases}$\\
$b$ is the number of points of discontinuity of $f_2(x)=\mathrm{sgn}\left((x-1)(x-2)^2(x-3)\right)$\\
$c$ is the number of points of non-differentiability of $f_3(x)=(x^2-1)\,|x^2-4x+3|+2(2-x)^{1/3}$\\
$d$ is the number of points of discontinuity of $f(x)=[2+3\sin x]\ \forall\ x\in(0,\pi)$\\
(where $[.]$ denotes greatest integer function)\\
On the basis of above information, answer the following questions :
\qn{9} Which of the following is/are true-
\begin{qoptions}
\opt{A}{$a^2+b^2=13$}
\opt{B}{$abcd=30$}
\opt{C}{$ac=8$}
\opt{D}{$a^2d=25$}
\end{qoptions}
\qn{10} Which of the following is/are true
\begin{qoptions}
\opt{A}{$a^2+2c+d=8$}
\opt{B}{$ab=8$}
\opt{C}{$a+2c=8$}
\opt{D}{$2a+3c+3d=25$}
\end{qoptions}
\sect{SECTION-II(i)}
\paraq{1}{2}
Let $f(x)=\tan\left(\dfrac{\pi(x+1)}{8}\right)+a$ \& $g(x)=x^2+bx+1$. If $h(x)=\displaystyle\lim_{n\to\infty}\dfrac{f(x)+(x^2)^ng(x)}{1+(x^2)^n}$ is continuous at $x=\pm1$ then
\qn{1} $a^2+b^2$ is equal to
\qn{2} minimum value of $g(x)$ is
\paraq{3}{4}
Let $f(x)=\cos^{-1}\left(\ell n\left(\dfrac{x+1}{x}\right)\right)$, then
\qn{3} The number of integers that does NOT lie in the Domain of $f(x)$, is
\qn{4} The number of integers lying in the range of $f(x)$, is
\paraq{5}{6}
Let $f(t)=\left(\sqrt t+\sqrt5\right)^2$ ($t\geqslant0$) and for sequence $\{T_n\}$, $T_1=5$, $n\geq2$, $T_n>0$, $S_n=f(S_{n-1})$, (where $T_n$ denotes $n^{\text{th}}$ term of the sequence and $S_n$ denotes sum of $n$ terms)\\
On the basis of above information, answer the following questions :
\qn{5} If $(T_{2018}+S_{25})$ is equal to $P$ then value of $\dfrac{P}{100}$ is-
\qn{6} If value of $n$ for which $\left(S_n-\dfrac73T_n\right)$ is minimum is $\dfrac ab$ ($a$ and $b$ are co-prime), then $a+b$ is-
\sect{SECTION-II(ii)}
\qn{7} Let $f(x)=ax^9+b\sin x+cx^2\,\mathrm{sgn}(x)+\dfrac{e^x-e^{-x}}{e^x+e^{-x}}$ be defined on set of real numbers (where $a>0$, $b,c\in R$). If $f(-5)=5$, $f(-2)=-3$, then the minimum number of zeroes of the equation $f(x)=0$, is :\\
(Where $\mathrm{sgn}(x)$ denotes signum function of $x$)
\qn{8} If $f(x)=\begin{cases}\dfrac{\sin^{-1}(\{x\}^2+\{x\}^3)-x}{\{x\}}, & x\neq0\\ -1, & x=0\end{cases}$, then the absolute value of RHD of $f(x)$ at $x=0$, is\\
(where $\{.\}$ denotes fractional part function)
\qn{9} Let $\{a_n\}$ be an arithmetic sequence with first term as 77 and common difference $-3$, then the maximum value of $\left(\displaystyle\sum_{k=1}^{n}a_k\right)$ is
\end{multicols}
\end{document}
